\documentclass[twocolumn]{aastex631}
\usepackage{amsmath}
\usepackage{graphicx}

\begin{document}

\title{Learning the Orbital-Uncertainty Propagation Operator: A Neural
Surrogate for Monte-Carlo Impact Monitoring of Near-Earth Asteroids}

\author{...}
\affiliation{...}

\begin{abstract}
Impact monitoring of near-Earth asteroids (NEAs) relies on propagating the
orbit-determination uncertainty into the future, traditionally by integrating
thousands of virtual asteroids (clones) sampled from the orbital covariance
with high-fidelity $N$-body dynamics (e.g.\ JPL Sentry). We present a neural
surrogate that learns this uncertainty-propagation operator directly: given the
cometary orbital elements and the $6\times6$ covariance matrix of an NEA, a
mixture-density network (MDN) deep ensemble predicts the full probability
distribution of the minimum Earth-approach distance in ten consecutive
decadal windows spanning 100~yr. Ground truth is generated with REBOUND
(IAS15) including all eight planets, the Moon, and general-relativistic
corrections, validated against JPL Horizons ephemerides. Trained on
$\sim$2{,}000 real NEAs from the JPL Small-Body Database with MOID
$<0.05$~au, the surrogate reproduces the Monte-Carlo close-approach
distributions with well-calibrated uncertainties (CRPS, PIT, and coverage
diagnostics) at a $\sim$$10^{4-5}\times$ speed-up over direct $N$-body clone
propagation, enabling population-scale risk triage for surveys such as
Rubin/LSST, which are expected to discover $\sim$$10^5$ new NEAs.
\end{abstract}

\keywords{Near-Earth objects --- Asteroids --- Impact monitoring --- Neural networks}

\section{Introduction}
% - Impact monitoring problem; Sentry / Sentry-II; clone MC cost
% - Deep learning in SSD: Hefele et al. 2020 (A&A 634, A45) binary classifier;
%   gap: no work learns the full uncertainty-propagation operator with UQ
% - LSST discovery rates motivate fast triage
TODO

\section{Data and Ground-Truth Generation}
\subsection{Target selection}
% SBDB query: 42,204 NEOs; condition code <= 3, data arc >= 365 d -> 10,907;
% Earth MOID < 0.05 au -> 3,424 targets with full 6x6 cometary covariance
TODO

\subsection{$N$-body Monte-Carlo propagation}
% REBOUND IAS15, Sun + 8 planets + Moon (Horizons initial conditions,
% epoch JD 2461000.5), REBOUNDx full GR (c in code units).
% Validation vs Horizons: Apophis 10-yr error 2.57e6 km (no GR) -> 4.7e4 km (GR);
% residual attributable to asteroid perturbers + Yarkovsky in JPL dynamics.
% Clones: multivariate normal in (e, q, tp, Omega, omega, i) at covariance epoch;
% train/val 96 clones, test 256 clones; sampling cadence 2.4/1.8 d with
% quadratic refinement of encounter minima; 100-yr horizon, 10 decadal windows.
TODO

\section{Neural Surrogate}
% Features: elements (angles as sin/cos, orbital phase) + signed-log Cholesky
% of covariance (32-dim); window embedding; residual MLP; 8-component MDN;
% deep ensemble over 5 seeds; NLL training.
TODO

\section{Results}
% CRPS / NLL / W1 vs climatology baseline; calibration (PIT, coverage);
% quantile agreement; example distributions; speed-up benchmark;
% case studies (Apophis, Bennu).
TODO

\section{Discussion and Conclusions}
TODO

\bibliography{refs}
\end{document}
